\documentclass[11pt]{article}

\usepackage[margin=1in]{geometry}
\usepackage{booktabs}
\usepackage{amsmath}
\usepackage{hyperref}

\title{Dual-Domain MMD Adaptation: Regularized Re-Run Update}
\author{}
\date{\today}

\begin{document}
\maketitle

\section{What changed}

The synthetic$\to$real adaptation direction (\texttt{syn\_source\_real\_target}) was re-run with
two additional regularization mechanisms enabled, keeping everything else (data, architecture,
epochs, layer, kernel) identical to the original run:

\begin{itemize}
    \item \textbf{Linear MMD-weight decay}: $\lambda_{DA}$ now decays linearly from $1.0$ to
    $0.1$ over the course of training, instead of staying constant at $0.8$ throughout.
    \item \textbf{Bandwidth freeze at epoch 5}: the RBF kernel's bandwidth stops updating after
    epoch 5 and is held fixed for the remainder of training, instead of continuing to track the
    batch statistics via exponential moving average for all 100 epochs.
\end{itemize}

\section{Result}

\begin{table}[h!]
    \centering
    \begin{tabular}{lrrrr}
        \toprule
        Model & Eval domain & mAP50 & mAP50-95 \\
        \midrule
        baseline (syn-pretrained) & syn (in-domain) & 0.758 & 0.508 \\
        baseline (real-pretrained) & real (in-domain) & 0.896 & 0.641 \\
        \midrule
        syn\_source\_real\_target (original) & real (target) & 0.867 & 0.599 \\
        syn\_source\_real\_target (original) & syn (source) & 0.568 & 0.337 \\
        \midrule
        \textbf{syn\_source\_real\_target\_reg (new)} & \textbf{real (target)} & \textbf{0.886} & \textbf{0.618} \\
        \textbf{syn\_source\_real\_target\_reg (new)} & \textbf{syn (source)} & \textbf{0.674} & \textbf{0.396} \\
        \bottomrule
    \end{tabular}
    \caption{Original vs.\ regularized re-run of the synthetic$\to$real direction, evaluated on
    the same stride-3 subsampled validation sets used for the pretrain baselines.}
\end{table}

The regularized re-run improves on both fronts simultaneously: target-domain (real) mAP50 rises
from 0.867 to 0.886, closing most of the remaining gap to the in-domain baseline (0.896). More
notably, source-domain (syn) performance -- the forgetting metric -- recovers substantially,
from 0.568 to 0.674 mAP50, more than halving the regression relative to the syn baseline (0.758).

\section{The EMA bandwidth mechanism}

The RBF kernel's bandwidth is not fixed; it is estimated as an exponential moving average of the
batch's mean pairwise squared distance, updated every training step. As the two domains' features
grow closer together under MMD alignment, this estimate shrinks too, which sharpens the kernel
and rewards the optimizer for collapsing the representation further rather than genuinely
matching the two distributions -- a feedback loop that self-reinforces as training progresses.
Freezing the bandwidth after a short warm-up period (epoch 5, by which point the EMA has already
smoothed over hundreds of batches) breaks this loop while still using a representative early-training
estimate, which is the most likely explanation for the reduced forgetting observed above.

\end{document}
